\section{形式积分}
\label{sec:mcalc-form-integration}

在一张坐标图中，$n$形式$f\,dx^1\wedge\cdots\wedge dx^n$的自然积分是$\int f\,\mathrm dx^1\cdots\mathrm dx^n$。问题出在换图：由式\eqref{eq:mcalc-top-form-change}，形式的系数按Jacobian行列式$\det\bm A$变换，而多元积分的换元公式乘的是$\lvert\det\bm A\rvert$。只有当所有换图的Jacobian行列式都为正时，两种变换才一致，积分才与坐标无关。这就要求先为流形规定一个“正方向”。本节先建立定向，再定义形式的积分；这里只需要光滑结构和定向，不需要黎曼度量。

\subsection{定向}
\label{sec:mcalc-orientation}

\paragraph{线性空间的定向}
平面上的有序基$(\bm e_1,\bm e_2)$与$(\bm e_2,\bm e_1)$描述了相反的旋转方向。一般地，$n$维线性空间$V$（$n\geq1$）的两组有序基之间的过渡矩阵若行列式为正，就称它们同向。同向是等价关系，恰把全部有序基分为两类；选定其中一类称为$V$的一个\term{定向}（orientation），这一类中的基称为正向基。零维空间约定定向就是选取符号$\pm1$。

\paragraph{流形的定向}
流形的每个切空间都可以选定向，问题是能否让这些选择随点连续地变化。用坐标图来表述：若一族坐标图覆盖$M$，并且任意两张图重叠处坐标变换的Jacobian行列式处处为正，称之为\term{定向图册}（oriented atlas）。

\begin{definition}[可定向流形]
\label{def:mcalc-orientable}
若光滑流形$M$存在定向图册，称$M$\term{可定向}（orientable）。选定一个定向图册（并约定与它的每张图同向的坐标图都算入），称为给$M$指定了一个定向；与之同向的坐标图称为正向坐标图，其坐标基在每点给出切空间的正向基。
\end{definition}

定向还有一种等价的刻画：$M$可定向当且仅当存在处处不为零的$n$形式。一个方向的推导很直接：若$\mu$处处非零，在每张连通的坐标图中$\mu=f\,dx^1\wedge\cdots\wedge dx^n$，$f$不变号，必要时把$x^1$换成$-x^1$可使$f>0$；这样调整后的图两两重叠处，由式\eqref{eq:mcalc-top-form-change}，Jacobian行列式等于两个正系数之比，为正。反方向用单位分解：给定定向图册与从属的单位分解$\{\psi_\alpha\}$，令$\mu=\sum_\alpha\psi_\alpha\,dx_\alpha^1\wedge\cdots\wedge dx_\alpha^n$；在任一正向坐标图中，每一项的系数是$\psi_\alpha$乘一个正的Jacobian行列式，非负且至少一项为正，所以$\mu$处处非零。

\begin{example}[可定向与不可定向]
\label{ex:mcalc-orientability}
$\mathbb R^n$及其开集可定向。单位球面$\mathbb S^2$可定向：$n$形式$\mu=x\,dy\wedge dz+y\,dz\wedge dx+z\,dx\wedge dy$限制到球面上处处非零，因为它在切向量$u,v$上的值是$\det[p,u,v]$，而$p$与$T_p\mathbb S^2$正交，当$u,v$线性无关时行列式非零。

M\"obius带是不可定向的典型。把正方形$[0,1]\times(-1,1)$的左右两边按$(0,s)\sim(1,-s)$粘合。沿中线$s=0$绕一圈，横向坐标$s$被翻转：任何试图沿中线连续地选取的正向基$(\partial_t,\partial_s)$，回到出发点时变成$(\partial_t,-\partial_s)$，定向反转。若存在定向图册，沿中线逐图传递定向就会产生矛盾，所以不存在定向图册。
\end{example}

\paragraph{边界的诱导定向}
Stokes公式要把区域上的积分与边界上的积分联系起来，必须先规定边界的定向。为使符号简洁，本章取带边界流形的局部模型为半空间$\mathbb H^n=\{x\in\mathbb R^n:x^1\leq0\}$，其边界为$\{x^1=0\}$；第~\ref{chap:manifold-basics}章若采用其他半空间，经坐标的置换与反射即可化为这一形式。在边界点$p$，称不切于边界、指向流形外侧的切向量为外法向，在这一模型中$\partial_1$指向外侧。

\begin{definition}[诱导定向]
\label{def:mcalc-boundary-orientation}
设$M$是定向的带边界流形，$n\geq2$，$p\in\partial M$，$N$是$p$处的外向切向量。称$T_p\partial M$的有序基$(v_1,\ldots,v_{n-1})$为正向基，当且仅当$(N,v_1,\ldots,v_{n-1})$是$T_pM$的正向基。这样得到的$\partial M$的定向称为\term{诱导定向}（induced orientation）。$n=1$时，边界点若外向切向量与$M$的正向一致取符号$+1$，否则取$-1$。
\end{definition}

“外法向在前”这条约定与正向的半空间模型相容：在$\mathbb H^n$中外向量$\partial_1$排在最前，$(\partial_1,\partial_2,\ldots,\partial_n)$是正向基，所以$\partial\mathbb H^n$以坐标$(x^2,\ldots,x^n)$的标准定向为诱导定向。两个低维例子可以检验它。区间$[a,b]$取通常定向时，右端点$b$处外向量与正向一致，得$+1$，左端点$a$处得$-1$。平面单位圆盘在$(1,0)$处的外法向为$\bm e_1$，$(\bm e_1,v)$正向要求$v$指向$\bm e_2$方向，所以边界圆取逆时针方向。

\subsection{微分形式的积分}
\label{sec:mcalc-integration-forms}

设$M$是定向的$n$维流形（可以带边界），$\omega$是$n$形式。积分的定义分两步：先处理支撑落在一张正向坐标图内的形式，再用单位分解拼接。这里要求$\omega$具有紧支撑，即$\operatorname{supp}\omega=\overline{\{p:\omega_p\neq0\}}$是紧集；这保证每一步都只涉及有限多项、有限的积分。

\paragraph{单张坐标图中的积分}
设$\operatorname{supp}\omega$含于正向坐标图$(U,\varphi)$的定义域内，在其中$\omega=f\,dx^1\wedge\cdots\wedge dx^n$。定义
\begin{equation}
\int_M\omega=\int_{\varphi(U)}f(x)\,\mathrm dx^1\cdots\mathrm dx^n .
\label{eq:mcalc-chart-integral}
\end{equation}
若$(V,\chi)$是另一张同样包含支撑的正向坐标图，坐标为$y$，由式\eqref{eq:mcalc-top-form-change}，$\omega=f\det\bm A\,dy^1\wedge\cdots\wedge dy^n$，其中$\bm A=[\partial x^i/\partial y^a]$且$\det\bm A>0$。多元积分的换元公式给出
\[
\int f(x)\,\mathrm dx=\int f(x(y))\,\lvert\det\bm A(y)\rvert\,\mathrm dy=\int f(x(y))\det\bm A(y)\,\mathrm dy ,
\]
最后一步正用到$\det\bm A>0$。所以两张图给出相同的值。换元公式可参考\BookChapterRef{chap:region-integrals}{卷二“区域积分”章}。

\paragraph{用单位分解拼接}
一般情形取正向坐标图覆盖$\{(U_\alpha,\varphi_\alpha)\}$与从属的单位分解$\{\psi_\alpha\}$。每个$\psi_\alpha\omega$的支撑含于$U_\alpha$，紧支撑与局部有限性保证只有有限个$\psi_\alpha\omega$不为零。定义
\begin{equation}
\int_M\omega=\sum_\alpha\int_M\psi_\alpha\,\omega ,
\label{eq:mcalc-form-integral}
\end{equation}
右端每项按式\eqref{eq:mcalc-chart-integral}计算。

\begin{theorem}[形式积分的良定性]
\label{thm:mcalc-form-integral-welldefined}
式\eqref{eq:mcalc-form-integral}的值不依赖正向坐标覆盖与单位分解的选取，对$\omega$线性，并且把$M$的定向换成相反定向时变号。
\end{theorem}

\begin{proof}
设另有正向覆盖$\{(V_\beta,\chi_\beta)\}$与单位分解$\{\rho_\beta\}$。每个$\psi_\alpha\rho_\beta\omega$的支撑同时含于$U_\alpha$与$V_\beta$，由单张图中积分与坐标无关，它在两张图中的积分相同。于是
\begin{equation*}
\begin{aligned}
\sum_\alpha\int\psi_\alpha\omega&=\sum_\alpha\int\Bigl(\sum_\beta\rho_\beta\Bigr)\psi_\alpha\omega=\sum_\alpha\sum_\beta\int\psi_\alpha\rho_\beta\omega\\
&=\sum_\beta\int\Bigl(\sum_\alpha\psi_\alpha\Bigr)\rho_\beta\omega=\sum_\beta\int\rho_\beta\omega ,
\end{aligned}
\end{equation*}
其中交换求和次序是允许的，因为只有有限个非零项；中间两步分别用$\sum_\beta\rho_\beta=1$与$\sum_\alpha\psi_\alpha=1$及单张图积分的线性。线性直接由定义得到。换成相反定向时，正向坐标图须换成例如把$x^1$换成$-x^1$的图，系数$f$变成$-f$，而积分区域的换元带来的是$\lvert\det\rvert=1$，所以每项变号。
\end{proof}

紧支撑条件也可以放宽：在每张正向坐标图中，若系数$f$关于Lebesgue测度绝对可积，并且$\sum_\alpha\int\lvert\psi_\alpha f_\alpha\rvert\,\mathrm dx<\infty$，同样的论证仍成立，此时交换求和次序由绝对收敛保证。

\paragraph{子流形上的积分}
要在曲线上积分$1$形式、在曲面上积分$2$形式，只需先拉回。设$S$是$k$维定向流形，$\iota:S\to M$是嵌入（或只是光滑映射），$\omega\in\Omega^k(M)$且$\iota^*\omega$紧支撑，定义
\[
\int_S\omega=\int_S\iota^*\omega .
\]
更一般地，若$F:M\to N$是两个定向$n$维流形之间的微分同胚，则对紧支撑的$n$形式$\omega$有$\int_MF^*\omega=\pm\int_N\omega$，$F$保持定向时取正号，反转定向时取负号。这正是坐标图中积分与坐标无关这一论证的整体形式。

\begin{example}[两个积分]
\label{ex:mcalc-form-integral-examples}
单位圆盘$D$取标准定向，求$\int_Ddx\wedge dy$。用极坐标$F(r,\theta)=(r\cos\theta,r\sin\theta)$，$(r,\theta)\in(0,1)\times(0,2\pi)$，它是到$D$去掉一条半径的保持定向的微分同胚（Jacobian行列式$r>0$），去掉的部分在积分中可忽略。由例~\ref{ex:mcalc-pullback-examples}，$F^*(dx\wedge dy)=r\,dr\wedge d\theta$，所以积分为$\int_0^{2\pi}\int_0^1r\,\mathrm dr\,\mathrm d\theta=\pi$，即圆盘的面积。

单位圆取逆时针定向，求$\int_{\mathbb S^1}(x\,dy-y\,dx)$。参数化$c(t)=(\cos t,\sin t)$，$t\in(0,2\pi)$，保持定向且只漏掉一点。由例~\ref{ex:mcalc-pullback-examples}，$c^*(x\,dy-y\,dx)=dt$，积分为$2\pi$。改取顺时针定向，结果变为$-2\pi$。上一节的闭形式$\omega=(x\,dy-y\,dx)/(x^2+y^2)$在单位圆上与这个形式相同，积分非零，正是它不恰当的另一种表现：下一节将看到，恰当形式在无边界的紧流形上积分为零。
\end{example}

例中“漏掉的部分可忽略”需要说明依据：去掉的是坐标图中的一个零测集，形式积分在单张图内就是Lebesgue积分，零测集不影响值。在后面的计算中，常用只漏掉零测集的参数化来代替单位分解。
